% -*- coding: utf-8 -*-
% !TEX program = xelatex
\documentclass[plain,noanswer,medmath]{../../template/jnuexam}
\usepackage{../../template/kysx}

\begin{document}


\examtitle{2018}{1}


\makepart{选择题}{1～8小题，每小题4分，共32分}


\begin{problem}
下列函数中，在$x = 0$处不可导的是\pickout{}
\begin{abcd}
\item $f(x) = |x|\sin |x|$．
\item $f(x) = |x|\sin \sqrt{|x|}$．
\item $f(x) = \cos |x|$．
\item $f(x) = \cos \sqrt{|x|}$．
\end{abcd}
\end{problem}


\begin{problem}
过点$(1,0,0)$，$(0,1,0)$，且与曲面$z = x^2 + y^2$相切的平面为\pickout{}
\begin{abcd}
\item $z = 0$与$x + y - z = 1$．
\item $z = 0$与$2x + 2y - z = 2$．
\item $y = x$与$x + y - z = 1$．
\item $y = x$与$2x + 2y - z = 2$．
\end{abcd}
\end{problem}


\begin{problem}
$\sum_{n = 0}^{\infty}(- 1)^n\frac{2n + 3}{\left(2n + 1 \right)!} =$\pickout{}
\begin{abcd}
\item $\sin 1 + \cos 1$．
\item $2\sin 1 + \cos 1$．
\item $2\sin 1 + 2\cos 1$．
\item $3\sin 1 + 2\cos 1$．
\end{abcd}
\end{problem}


\begin{problem}
设$M = \int_{-\frac{\pi}{2}}^{\frac{\pi}{2}} \frac{\left(1 + x \right)^2}{1 + x^2}\dx$,
  $N = \int_{-\frac{\pi}{2}}^{\frac{\pi}{2}} \frac{1 + x}{\e^x}\dx$,
  $K = \int_{-\frac{\pi}{2}}^{\frac{\pi}{2}} \left(1 + \sqrt{\cos x}\right)\dx$，
则\pickout{}
\begin{abcd}
\item $M > N > K$．
\item $M > K > N$．
\item $K > M > N$．
\item $K > N > M$．
\end{abcd}
\end{problem}


\begin{problem}
下列矩阵中，与矩阵$\begin{pmatrix}
  1 & 1 & 0 \\
  0 & 1 & 1 \\
  0 & 0 & 1
\end{pmatrix}$相似的为\pickout{}
\begin{abcd}
\item $\begin{pmatrix}
  1 & 1 & -1 \\
  0 & 1 & 1 \\
  0 & 0 & 1
\end{pmatrix}$．
\item $\begin{pmatrix}
  1 & 0 & -1 \\
  0 & 1 & 1 \\
  0 & 0 & 1
\end{pmatrix}$．
\item $\begin{pmatrix}
  1 & 1 & -1 \\
  0 & 1 & 0 \\
  0 & 0 & 1
\end{pmatrix}$．
\item $\begin{pmatrix}
  1 & 0 & -1 \\
  0 & 1 & 0 \\
  0 & 0 & 1
\end{pmatrix}$．
\end{abcd}
\end{problem}


\begin{problem}
设$A,B$为$n$阶矩阵，记$r(X)$为矩阵$X$的秩，$(X,Y)$表示分块矩阵，则\pickout{}
\begin{abcd}
\item $r(A,AB) = r(A)$．
\item $r(A,BA) = r(A)$．
\item $r(A,B) = \max\{r(A),r(B)\}$．
\item $r(A,B) = r(A^T,B^T)$．
\end{abcd}
\end{problem}


\begin{problem}
设$f(x)$为某分布的概率密度函数，$f\left(1 + x \right) = f(1 - x)$，$\int_0^2 f(x)\dx = 0.6$，
则$P\{X < 0\} =$\pickout{}
\begin{abcd}
\item $0.2$．
\item $0.3$．
\item $0.4$．
\item $0.6$．
\end{abcd}
\end{problem}


\begin{problem}
给定总体$X\sim N\left(\mu ,\sigma^2 \right)$，$\sigma^2$已知，给定样本$X_1,X_2, \cdots ,X_n$，
对总体均值$\mu$进行检验，令$H_0:\mu = \mu_0$，$H_1:\mu \ne \mu_0$，则\pickout{}
\begin{abcd}
\item 若显著性水$\alpha = 0.05$时拒绝$H_0$，则$\alpha = 0.01$时也拒绝$H_0$．
\item 若显著性水$\alpha = 0.05$时接受$H_0$，则$\alpha = 0.01$时拒绝$H_0$．
\item 若显著性水$\alpha = 0.05$时拒绝$H_0$，则$\alpha = 0.01$时接受$H_0$．
\item 若显著性水$\alpha = 0.05$时接受$H_0$，则$\alpha = 0.01$时也接受$H_0$．
\end{abcd}
\end{problem}


\makepart{填空题}{9～14小题，每小题4分，共24分}


\begin{problem}
若$\lim_{x \to 0} \left(\frac{1 - \tan x}{1 + \tan x} \right)^{\frac{1}{\sin kx}} = \e$，
则$k =$ \fillin{}．
\end{problem}


\begin{problem}
设函数$f(x)$具有$2$阶连续导数，若曲线$y = f(x)$过点$(0,0)$
且与曲线$y = 2^x$在点$(1,2)$处相切，则$\int_0^1 xf''(x)\dx =$ \fillin{}．
\end{problem}


\begin{problem}
设$F(x,y,z) = xy\vi - yz\vj + zx\vk$，则$\operatorname{rot}\vec{F}(1,1,0) =$ \fillin{}．
\end{problem}


\begin{problem}
设$L$为球面$x^2 + y^2 + z^2 = 1$与平面$x + y + z = 0$的交线，则$\oint_L xy\ds =$ \fillin{}．
\end{problem}


\begin{problem}
设$2$阶矩阵$A$有两个不同特征值，$\alpha_1,\alpha_2$是$A$的线性无关的特征向量，
且满足$A^2(\alpha_1 + \alpha_2) = (\alpha_1 + \alpha_2)$，
则$|A| =$ \fillin{}．
\end{problem}


\begin{problem}
设随机事件$A$与$B$相互独立，$A$与$C$相互独立，$BC = \emptyset$，若
$$P(A) = P(B) = \frac{1}{2},\quad P\left(AC|AB \cup C \right) = \frac{1}{4},$$
则$P(C) =$ \fillin{}．
\end{problem}


\makepart{解答题}{15～23小题，共94分}


\begin{problem}[points=10]
求不定积分$\int \e^{2x}\arctan \sqrt{\e^x - 1} \dx$．
\end{problem}


\begin{problem}[points=10]
将长为$2m$的铁丝分成三段，依次围成圆、正方形与正三角形，
三个图形的面积之和是否存在最小值？若存在，求出最小值．
\end{problem}


\begin{problem}[points=10]
设$\Sigma$是曲面$x = \sqrt{1 - 3y^2 - 3z^2}$的前侧，计算曲面积分
$$\iint_{\Sigma}x\dy\dz + \left(y^3 + z \right)\dx\dz + z^3 \dx\dy.$$
\end{problem}


\begin{problem}[points=10]
已知微分方程$y' + y = f(x)$，其中$f(x)$是$R$上的连续函数．
\begin{enumerate}
\item 若$f(x) = x$时，求微分方程的通解；
\item 若$f(x)$周期为$T$函数时，证明：方程存在唯一的以$T$为周期的解．
\end{enumerate}
\end{problem}


\begin{problem}[points=10]
设数列$\left\{x_n \right\}$满足：$x_1 > 0$，$x_n \e^{x_{n + 1}} = \e^{x_n} - 1$（$n = 1,2, \cdots$）．
证明$\left\{x_n \right\}$收敛，并求$\lim_{n \to \infty} x_n$．
\end{problem}


\begin{problem}[points=11]
设实二次型$f(x_1,x_2,x_3) = (x_1 - x_2 + x_3)^2 + (x_2 + x_3)^2 + (x_1 + ax_3)^2$，其中$a$是参数．\par
\begin{enumerate*}
\item 求$f(x_1,x_2,x_3) = 0$的解；
\item 求$f(x_1,x_2,x_3)$的规范形．
\end{enumerate*}
\end{problem}


\begin{problem}[points=11]
已知$a$是常数，且矩阵$A = \begin{pmatrix}
  1 & 2 & a \\
  1 & 3 & 0 \\
  2 & 7 & -a
\end{pmatrix}$可经初等列变换化为矩阵$B = \begin{pmatrix}
   1 & a & 2 \\
   0 & 1 & 1 \\
  -1 & 1 & 1
\end{pmatrix}$．\par
\begin{enumerate*}
\item 求$a$；
\item 求满足$AP = B$的可逆矩阵$P$．
\end{enumerate*}
\end{problem}


\begin{problem}[points=11]
已知随机变量$X$，$Y$相互独立，且$P\{X = 1\} = P\{X = - 1\} = \frac{1}{2}$，
$Y$服从参数为$\lambda$的泊松分布，令$Z = XY$．\par
\begin{enumerate*}
\item 求$\Cov \left(X,Z \right)$；
\item 求$Z$的概率分布．
\end{enumerate*}
\end{problem}


\begin{problem}[points=11]
设$X$的概率密度为
$$f\left(x,\sigma \right) = \frac{1}{2\sigma}\e^{- \frac{|x|}{\sigma}},\quad -\infty < x < +\infty ,$$
其中$\sigma \in (0,+\infty)$为未知参数，$X_1,X_2, \cdots X_n$为来自总体$X$的简单随机样本，
记$\sigma$的最大似然估计量为$\hat{\sigma}$．\par
\begin{enumerate*}
\item 求$\hat{\sigma}$；
\item 求$E\left(\hat{\sigma} \right)$，$D\left(\hat{\sigma} \right)$．
\end{enumerate*}
\end{problem}


\end{document}
